\documentclass[11pt]{article}
\usepackage[margin=1in]{geometry}
\usepackage{booktabs}
\usepackage{amsmath}
\usepackage{longtable}
\title{Project 3: Water Resources Optimization for Reservoir Dispatch}
\author{Huang Qiwei \\ 3125301141 \\ Xi'an Jiaotong University}
\date{Software Development 2026}

\begin{document}
\maketitle

\section*{Abstract}
This experiment solves a seven-day reservoir dispatch problem under hydropower, ecological release, storage, and mass-balance constraints. The implementation prioritizes SciPy's SLSQP constrained optimizer when the dependency is installed. A linear fallback solver is included only to keep the workflow reproducible in lightweight environments where SciPy or Matplotlib may be unavailable.

\section*{Problem Definition}
The reservoir starts with storage $V_0=500{,}000$ m$^3$. Daily release decisions must respect storage bounds, release bounds, and the water-balance equation:
\[
V_{t+1}=V_t+(I_t-Q_t)\Delta t.
\]
The ecological hard-constraint case requires $Q_t \ge 10$ m$^3$/s. The comparison case allows lower releases so the cost of ecological protection can be estimated.

\section*{Objectives and Constraints}
\begin{longtable}{ll}
\toprule
Item & Definition \\
\midrule
Decision variables & Seven daily releases $Q_1,\ldots,Q_7$ \\
Revenue objective & Maximize hydropower revenue weighted by daily price \\
Ecological metric & Sum of deficits below 10 m$^3$/s \\
Storage bounds & $100{,}000 \le V_t \le 1{,}000{,}000$ m$^3$ \\
Release bounds & $0 \le Q_t \le 100$ m$^3$/s, with a hard ecology case of $Q_t \ge 10$ \\
Mass balance & Daily storage update from inflow and release \\
\bottomrule
\end{longtable}

\section*{Implementation}
The main script is \texttt{reservoir\_optimization.py}. When SciPy is installed from \texttt{requirements.txt}, the model uses \texttt{scipy.optimize.minimize} with SLSQP and nonlinear storage constraints. This is the intended strict optimization route. If SciPy is unavailable, the script reports that condition and uses a linear fallback to generate feasible schedules and figures for review, without pretending that SLSQP was run.

\section*{Trade-off Analysis}
The script compares a hard ecological release solution with a revenue-prioritized comparison solution. It then computes a Pareto-style trade-off figure, \texttt{tradeoff\_analysis.png}, showing the revenue change associated with ecological deficit. This supports interpretation of how much revenue is sacrificed to maintain downstream flow.

\section*{Results}
The exported \texttt{optimal\_schedule.csv} contains the daily inflow, release, storage, price, and revenue. The \texttt{validation\_report.txt} confirms storage bounds, release bounds, and mass balance. In the available run, the hard ecological solution satisfies all constraints and provides a defensible reservoir operating schedule.

\section*{Validation}
Validation is central because optimization code can return numerically plausible but physically invalid schedules. The report checks:
\begin{itemize}
  \item all releases stay within the relevant bounds,
  \item storage never falls below minimum or exceeds maximum,
  \item the daily mass balance closes within tolerance,
  \item total revenue is calculated consistently from releases and prices.
\end{itemize}

\section*{AI-assisted Development Notes}
AI assistance was used to formulate the objective, implement the optimizer, design constraint checks, and produce the trade-off visualization. Each generated step was checked against the reservoir mass-balance equation and the experiment statement.

\section*{Limitations and Future Work}
The fallback solver is a reproducibility aid, not a replacement for the SciPy SLSQP solution expected in a strict optimization environment. Future extensions could include inflow uncertainty, rolling-horizon model predictive control, water quality constraints, and multi-objective algorithms beyond weighted-sum analysis.

\end{document}
